\documentclass[UTF8,11pt,a4paper]{ctexart}

\usepackage{geometry}
\geometry{margin=2.2cm}
\usepackage{amsmath,amssymb}
\usepackage{booktabs}
\usepackage{hyperref}
\usepackage{enumitem}
\usepackage{xcolor}
\usepackage{listings}
\usepackage{longtable}
\usepackage{fancyhdr}

\definecolor{codebg}{RGB}{248,248,248}
\definecolor{codeframe}{RGB}{210,210,210}
\lstset{
  basicstyle=\ttfamily\small,
  backgroundcolor=\color{codebg},
  frame=single,
  rulecolor=\color{codeframe},
  breaklines=true,
  columns=fullflexible,
  keepspaces=true,
  tabsize=2,
  showstringspaces=false
}

\setlength{\headheight}{14pt}
\pagestyle{fancy}
\fancyhf{}
\lhead{QUANTUM 放电类粒子管理方案}
\rhead{merge-only 主导版}
\cfoot{\thepage}

\title{QUANTUM：放电类算例粒子管理提升方案\\
{\large 以粒子合并为主，原则上不做常规分裂}}
\author{}
\date{2026年4月13日}

\begin{document}
\maketitle

\begin{abstract}
本文针对 QUANTUM 当前粒子管理策略在放电类算例中的两个核心问题提出改进方案：
其一，频繁 merge/split 会引入明显的统计波动，尤其在 sheath、近壁面和低密度区更容易放大噪声；
其二，现有粒子管理本身需要额外的分箱、复制和重建，某些情况下甚至比完全不做管理更慢。

基于对当前代码路径的审计，本文建议将粒子管理的目标从“同时维持上下限”改为“只压制上尾，不主动补低尾”，也就是：
以 merge 为主，split 默认关闭，仅保留非常罕见的应急兜底；
同时把主要控制点前移到 spawn 端和总量上限，而不是在全域内频繁做 per-cell 粒子数修补。

该方案既符合放电模拟对低噪声和稳定电流统计的要求，也更贴合当前 QUANTUM 的工程现状，可作为显式 PIC 的直接改造方向，并可在后续扩展到隐式 PIC。
\end{abstract}

\tableofcontents

\section{结论先行}

对于放电类算例，建议采用如下总原则：
\begin{enumerate}[itemsep=3pt]
  \item \textbf{只控制宏粒子数的上尾，不填补下尾。} 也就是保留 merge，取消常规 per-cell minimum 驱动的 split。
  \item \textbf{把控制重点放在 spawn 入口和总量上限。} 先防止新生粒子把总量推高，再做低频、温和的 bulk merge。
  \item \textbf{默认不在 sheath、cut-cell、近电极区积极做粒子管理。} 这些区域物理梯度大，统计扰动最容易传导到电流和壁面通量。
  \item \textbf{如果确实需要 split，只允许极低频、极小规模、二分式应急 split。} 不再使用当前的常规四分裂作为生产默认策略。
\end{enumerate}

换句话说，这份方案的核心不是“把现有 merge/split 调得更保守”，而是直接改变策略目标：
从“维持某个 per-cell 粒子数窗口”改为“在不明显破坏统计性质的前提下，限制粒子数上尾和权重上尾”。

\section{对当前实现的审计结论}

\subsection{当前代码已经具备什么}
从现有代码看，粒子管理相关能力主要分成三层：
\begin{itemize}[itemsep=3pt]
  \item 配置结构在 \texttt{ParticleManagementConfig.h} 中，已经支持全局 soft/hard cap、per-cell 上下限、merge、split、spawn thinning 和 report。
  \item 解析入口在 \texttt{DictionaryPIC.cpp} 中，也就是说 \texttt{pic.dict} 里的 \texttt{ParticleManagement} 块已经可以被显式和隐式共用配置层读入。
  \item 实际算法在 \texttt{ParticleManager.h} 中，当前可见的执行调用点在 \texttt{ExplicitPICModule.cpp} 的 spawn commit 之后。
\end{itemize}

这意味着：\textbf{当前已经真正执行粒子管理的是显式 PIC 路径。}
就本次审计范围内看到的代码而言，隐式 PIC 侧虽然能读到同名配置，但没有看到等价的 \texttt{ParticleManager::Apply(...)} 调用链。
因此，如果某个隐式算例写了 \texttt{ParticleManagement} 配置块，现阶段更可能是“配置存在但未真正生效”。

\subsection{当前实现为什么容易引起波动和额外开销}

当前实现的主流程是：
\begin{enumerate}[itemsep=3pt]
  \item 对某一物种按 cell 重新分箱；
  \item 当 cell 内宏粒子数高于 \texttt{perCellMax} 时做 merge；
  \item 当 cell 内宏粒子数低于 \texttt{perCellMin} 或单粒子权重大于 \texttt{splitWMax} 时做 split；
  \item 生成一个新的粒子数组并整体替换旧容器。
\end{enumerate}

这里有四个与放电问题高度相关的缺点：

\paragraph{缺点 A：当前 merge 选对方式过于粗糙}
当前 \texttt{ParticleManager} 并不是在相空间中寻找最近邻，也不是按速度壳层或能量壳层做保守聚类；
它是把 cell 内粒子拷到局部数组后，直接取末尾两个粒子尝试配对合并。
虽然代码会用 \texttt{mergeMaxDeltaE} 做一次能量差阈值判断，但“候选对”的生成本身并不物理。
这会带来两个后果：
\begin{itemize}[itemsep=2pt]
  \item merge 结果对粒子排列顺序敏感；
  \item 即使通过了阈值，也不保证这就是局部最合适的配对。
\end{itemize}

\paragraph{缺点 B：当前 split 触发条件天然会把低密度区推向高波动}
当前 split 的两个触发条件分别是：
\begin{itemize}[itemsep=2pt]
  \item cell 内粒子数低于 \texttt{perCellMin}；
  \item 单粒子权重大于 \texttt{splitWMax}。
\end{itemize}
对放电算例而言，第一个条件问题最大。
sheath、近壁面、边界层和一些耗尽区本来就会出现低粒子数，这是物理解的结果，不一定说明采样已经失效。
如果强行用 split 把这些区域“补到某个最低粒子数”，就会把原本只是低样本的问题，变成主动注入更多相关粒子的问题。

\paragraph{缺点 C：当前 split 方式过于激进}
当前 split 不是一分为二，而是把一个粒子拆成四个位置抖动的拷贝，速度基本保持不变，权重变成四分之一。
这意味着一次 split 会让局部宏粒子数瞬间增加 3 个。
如果这种操作在低密度区、强场区或每隔很短步长反复触发，那么 CPU 开销和统计相关性都会变差。

\paragraph{缺点 D：当前 PM 的成本本身不低}
每次 PM 执行都需要：
\begin{itemize}[itemsep=2pt]
  \item 为所有 cell 构建 bins；
  \item 将 cell 内粒子复制到临时容器；
  \item 构建 \texttt{newParticles} 并整体换回；
  \item 最后还要做一次嵌入边界内粒子清理。
\end{itemize}
如果 PM 频率较高，而 split 又不断增加粒子数，那么“粒子管理本身成为热点”是完全可能的。
这正好解释了为什么某些 case 会出现“开了 PM 反而比不开更慢”的体验。

\section{适合放电算例的新策略：只压上尾，不补下尾}

\subsection{目标函数的改变}
当前思路本质上是在控制每个 cell 的粒子数落在一个窗口内：
\[
N_{\min} \le N_{cell} \le N_{\max}.
\]
对于均匀、稳态、统计量以体平均为主的问题，这种思路有一定合理性。
但对放电类算例，更合适的目标是：
\[
\text{限制总粒子数上尾} + \text{限制局部极端过密 cell} + \text{限制权重重尾},
\]
而不是强行要求所有 cell 都达到某个最低粒子数。

因此推荐把控制目标改成三条：
\begin{enumerate}[itemsep=3pt]
  \item 总粒子数不能失控增长；
  \item bulk 区域不能出现极端过密单元；
  \item 不能形成过重的权重长尾，避免少数大权重粒子主导统计噪声。
\end{enumerate}

\subsection{生产默认策略 M1：merge-only 主导}

建议将放电类算例的生产默认模式定义为 M1：
\begin{itemize}[itemsep=3pt]
  \item \texttt{mergeEnabled = true}
  \item \texttt{splitEnabled = false}
  \item \texttt{perCellMin = 0}
  \item \texttt{thinningEnabled = true}
  \item \texttt{every = 20--100}，而不是每步都做
\end{itemize}

该模式的含义非常明确：
\begin{itemize}[itemsep=2pt]
  \item 允许低密度区保持低粒子数；
  \item 只在 cell 过密时做合并；
  \item 新生粒子过多时优先在 spawn 端做概率保留；
  \item 不再常规使用 split 去“修补”低密度区。
\end{itemize}

这比单纯把 \texttt{splitWMax} 调大更干净，因为它直接从策略层面承认：
\textbf{放电问题里最重要的是别把下尾当成错误。}

\subsection{生产默认策略 M1 的三层控制}

\paragraph{层 1：spawn 端保险丝}
对 MCC、SEE 和 source injection 产生的新粒子，先做总量视角的 thinning。
这一步的优点是它发生在 commit 之前，能直接抑制峰值内存和突发增长。

建议保留当前的概率保留思想：
\[
p = \min\left(1, \frac{N_{soft}}{N_{current}+N_{spawn}}\right),
\]
必要时再叠加 \texttt{hardCapTotal} 限制。

但对放电场景，建议新增两条语义约束：
\begin{enumerate}[itemsep=2pt]
  \item thinning 主要针对电子和碰撞高产物种，慢离子和关键诊断物种可单独关闭；
  \item 如果某次 thinning 会把粒子权重抬得过高，应优先降低执行频率或上调 source macro weight，而不是事后再靠 split 去补救。
\end{enumerate}

\paragraph{层 2：bulk 区域低频 merge}
merge 的目标不是“把所有 cell 压到同一个数”，而是只处理明显过密的 cell。
建议将 merge 触发改成带迟滞的双阈值：
\[
N_{cell} > N_{merge,on} \Rightarrow \text{允许 merge}, \qquad
N_{cell} < N_{merge,off} \Rightarrow \text{停止 merge},
\]
其中 $N_{merge,off} < N_{merge,on}$，避免同一 cell 在相邻 PM 周期里反复进出状态。

一个可行的经验值是：
\begin{itemize}[itemsep=2pt]
  \item 电子：$N_{merge,on}=96$，$N_{merge,off}=64$
  \item 离子：$N_{merge,on}=64$，$N_{merge,off}=40$
\end{itemize}
这不是固定真值，但比单一的 \texttt{perCellMax/perCellMin} 窗口更接近实际需求。

\paragraph{层 3：权重重尾监控，而不是常规 split}
大权重粒子确实会带来噪声，但并不意味着常规 split 是最佳解。
更推荐的做法是：
\begin{itemize}[itemsep=2pt]
  \item 统计每个物种的权重分位数，例如 p50、p90、p99、max；
  \item 统计 top 1\% 权重粒子的物理粒子份额；
  \item 当权重分布开始明显重尾时，优先回看 thinning 概率、spawn 频率和 merge 阈值；
  \item 只有在极端异常时才启用应急 split。
\end{itemize}

\section{建议替换的算法细节}

\subsection{把当前“尾部两两配对”改成更稳健的 merge 配对}
如果继续沿用当前 \texttt{ParticleManager} 框架，最值得先改的一点不是 split，而是 merge 配对策略。

建议的最小改造版如下：
\begin{enumerate}[itemsep=3pt]
  \item 仍然按 cell 分箱；
  \item 先按电荷符号和主速度方向分组，例如沿主电场方向的速度正负号分开；
  \item 在每个子组内按动能或速度模长排序；
  \item 只配对相邻项，且仍保留 \texttt{mergeMaxDeltaE} 阈值。
\end{enumerate}

这样做的好处是：
\begin{itemize}[itemsep=2pt]
  \item 不需要昂贵的全局最近邻搜索；
  \item 相比当前“容器尾部二元组”更稳定；
  \item 更适合放电问题里常见的各向异性速度分布。
\end{itemize}

如果后续还要继续提升，可以再发展到真正的 velocity-space binning 或 k-nearest merge，但那应该是第二阶段工作，不必一开始就做重型算法。

\subsection{取消常规 per-cell minimum}
建议从配置和实现两个层面都弱化甚至移除 \texttt{perCellMin}：
\begin{itemize}[itemsep=3pt]
  \item 在生产默认配置里直接设为 0；
  \item 在代码层把“低于最小粒子数就 split”从默认主路径里拿掉；
  \item 若保留该参数，也应仅作为实验性或诊断性开关存在。
\end{itemize}

这一步会直接消除一类最容易制造波动的 split 触发器。

\subsection{如果必须 split，改成低频二分式应急 split}
如果后续确实发现某些 case 存在非常少数的超大权重粒子，建议把 split 改为如下语义：
\begin{enumerate}[itemsep=3pt]
  \item 默认关闭；
  \item 只在权重超过物种目标权重若干倍时触发；
  \item 只允许一分为二，不做一分为四；
  \item 只在非 cut-cell、非近壁面、且本 cell 已经有足够邻居样本时触发；
  \item 同一 cell 或同一粒子需要有 cooldown，不能连续触发。
\end{enumerate}

这样的 split 更像故障保险，而不是统计维持手段。

\section{区域感知：比“全域统一 PM”更适合放电}

如果希望方案更稳，我建议加入区域感知，而不是在所有 cell 使用同一套 PM 参数。

\subsection{建议默认跳过的区域}
以下区域默认不做或少做 PM：
\begin{itemize}[itemsep=3pt]
  \item cut-cell 和紧邻 embedded boundary 的 cell；
  \item 紧邻电极、发射边界、入口边界的 cell；
  \item 电场极强或电势梯度极大的 cell；
  \item 电子能量分布明显非局域的区域。
\end{itemize}

原因很直接：这些地方的粒子分布最不接近“局部统计平稳”，也是 PM 最容易污染物理信号的地方。

\subsection{建议优先做 PM 的区域}
更适合 merge-only 的区域是：
\begin{itemize}[itemsep=3pt]
  \item 近似准中性的 bulk 区；
  \item 粒子数非常高但分布相对平滑的体区；
  \item 远离壁面、远离 cut-cell 的普通流通区。
\end{itemize}

这会把 PM 从“全局修补工具”变成“bulk 区压缩工具”，更符合放电问题的物理直觉。

\section{建议的配置接口}

\subsection{短期可直接落地的配置语义}
在不大改配置解析器的前提下，建议先把放电默认配置写成下面这种风格：

\begin{lstlisting}
ParticleManagement
{
    enabled true
    every 50

    softCapTotal 2000000
    hardCapTotal 2600000

    perCellMax 96
    perCellMin 0

    merge true
    mergeMaxDeltaEBySpecies (2.0 eV 0.5 eV 0.5 eV)
    mergeMaxPairsPerCell 16

    split false
    splitWMax 0
    splitMaxPerCell 0

    thinningEnabled true
    thinningMccSpawn true
    thinningSeeSpawn true
    minKeepProb 0.2

    reportEnabled true
    reportEvery 50
}
\end{lstlisting}

这组参数表达的含义是：
\begin{itemize}[itemsep=3pt]
  \item PM 低频执行；
  \item 只保留上限控制；
  \item split 明确关闭；
  \item spawn thinning 保留，但不允许保留概率太小，以免权重暴涨。
\end{itemize}

\subsection{中期建议增加的配置项}
如果下一步准备正式把该方案做成放电专用模式，建议在配置层新增这些字段：
\begin{itemize}[itemsep=3pt]
  \item \texttt{mergeOffThreshold}：配合 \texttt{perCellMax} 形成迟滞；
  \item \texttt{skipCutCell}：cut-cell 默认跳过 PM；
  \item \texttt{skipNearBoundaryLayers}：离壁面若干层单元内不做 PM；
  \item \texttt{splitMode = off | emergency}：明确 split 是不是只作兜底；
  \item \texttt{weightTailWarnP99}：权重分布告警阈值；
  \item \texttt{targetSpecies}：只对选定物种做 PM。
\end{itemize}

其中最关键的两个是 \texttt{mergeOffThreshold} 和 \texttt{targetSpecies}。
它们能让 PM 从“粗暴地全域启用”变成“按需、低频、局部启用”。

\section{分阶段落地路线}

\subsection{阶段 A：立刻可做的最小改动}
这个阶段不追求算法升级，只追求先把方向纠正：
\begin{enumerate}[itemsep=3pt]
  \item 放电 case 默认设 \texttt{splitEnabled=false}；
  \item 放电 case 默认设 \texttt{perCellMin=0}；
  \item PM cadence 改为低频，例如每 20 到 100 步一次；
  \item 保留 spawn thinning；
  \item 保留现有 report，但新增权重分位数输出。
\end{enumerate}

如果只做这一步，通常就已经能显著降低“PM 比不开还慢”和“PM 触发额外涨落”的问题。

\subsection{阶段 B：改造 merge 配对策略}
这个阶段建议替换当前尾部二元组配对，至少改成“分组后排序再邻接配对”。
它的代码侵入度不大，但对稳定性改善会很明显。

\subsection{阶段 C：加入区域感知}
此阶段的目标是让 PM 主要服务于 bulk，而不是干扰 sheath 与近壁面。
如果只能做一个区域规则，我建议优先加“cut-cell 与其邻域默认跳过 PM”。

\subsection{阶段 D：仅在必要时加入应急 split}
只有在观察到非常明显的权重重尾，并且该重尾已经实质性污染统计量时，才建议加入 emergency split。
而且应急 split 应当：
\begin{itemize}[itemsep=2pt]
  \item 默认关闭；
  \item 一分为二；
  \item 仅针对指定物种；
  \item 有 cooldown；
  \item 不在 cut-cell 和近壁面触发。
\end{itemize}

\subsection{阶段 E：扩展到隐式 PIC}
由于当前显式模块才是实际调用 PM 的主路径，因此如果后续希望让隐式算例真正受益，需要额外完成两件事：
\begin{enumerate}[itemsep=3pt]
  \item 在隐式 PIC 时间步中明确选择 PM 插入点；
  \item 复核 implicit orbit-averaged 权重、极化泊松和输出统计链路，确认 variable weight 或 spawn thinning 不会破坏现有假设。
\end{enumerate}

也就是说，\textbf{隐式侧不应只复制配置，而应明确复制语义和调用时机。}

\section{验收指标}

如果要判断这套新策略是否成功，建议不要只看总粒子数，还要同时跟踪以下指标：

\begin{longtable}{p{0.26\linewidth}p{0.68\linewidth}}
\toprule
指标 & 建议判据 \\
\midrule
总宏粒子数 & 不再出现持续失控增长，且 PM 打开后不应明显高于无 PM 基线的 CPU 时间 \\
电流与壁面通量 & 与无 PM 基线相比，平均值接近，短时波动不出现系统性放大 \\
电子密度与势分布 & bulk 区统计量保持一致，sheath 宽度与近壁电势不因 PM 出现异常台阶 \\
权重分布 & p99 与 max 不应长期漂移到极端大值；top 1\% 权重粒子的物理份额应受控 \\
PM 自身耗时 & PM 耗时占比应显著低于推进、沉积、泊松等主热点，至少不能超过不开 PM 时的总耗时收益 \\
\bottomrule
\end{longtable}

验收时建议始终保留一组“完全不开 PM”的对照算例。
对放电问题而言，PM 的目标不是让统计量看起来更平滑，而是在不明显改变统计规律的前提下控制成本和粒子数。

\section{推荐结论}

综合当前代码现状和放电问题的物理需求，我的建议是：
\begin{enumerate}[itemsep=3pt]
  \item \textbf{立刻停止把 split 当作常规手段。} 生产默认应为 merge-only 主导。
  \item \textbf{立刻取消 per-cell minimum 驱动。} 低密度区不是错误，不应默认补粒子。
  \item \textbf{优先保留 spawn thinning。} 它对峰值粒子数和内存更有效，也更符合当前工程结构。
  \item \textbf{优先改 merge 配对，而不是优先发明更复杂的 split。}
  \item \textbf{如果后续还觉得噪声来自大权重粒子，先查权重重尾来源，再决定是否需要 emergency split。}
\end{enumerate}

一句话概括就是：
\begin{center}
\large
放电类算例的粒子管理，应该是“少动粒子、只压上尾、尽量不补下尾”。
\end{center}

这比“merge/split 都开着，通过参数慢慢调”更符合当前 QUANTUM 的实现条件，也更符合工程上可控、可验证、可逐步落地的方向。

\end{document}